\section{Engineering, availability, and sustainability}
HCMO is released not as a single file but as a \textbf{tool-consumable package} designed so that downstream tools can resolve every component programmatically.

\textbf{Release manifest as a stable contract.} A machine-readable manifest (\nolinkurl{hcmo.yaml}) declares the ontology's identity (name, version, namespace, prefix, ontology and version IRIs), the ordered list of source modules, the generated distributions, the SHACL shapes, the competency-question index, and the example datasets. Its \emph{shape} is treated as an API and kept stable across releases, so a consumer (e.g. a synchronisation or question-answering layer) can discover paths from the manifest rather than hard-coding them.

\textbf{Modular sources and reproducible distributions.} The ontology is authored as five modular Turtle sources (\emph{core/bio/env/obs/tech}) plus a migration-only compatibility module. A build step merges them into a canonical, sorted serialisation, producing byte-identical output on re-run so that version-control diffs stay clean. The merged graph is published in three syntaxes --- Turtle, RDF/XML (OWL), and JSON-LD --- together with a flat term inventory (\nolinkurl{profile.json}: IRIs, labels, comments, and counts) intended for sync layers and user interfaces. Everything under the distribution directory is generated; only the modules are edited by hand. The manifest still identifies version 0.2.0, but the paper-matching graph contains 24 later commits, including the ISA/STATO, SemTS/SOSA, temporal-state, and QUDT work. A new version and tag must therefore be assigned before submission; the tagged 0.2.0 artifact is not presented as the evaluated paper release.

\textbf{Continuous validation.} A validation step parses every Turtle file, runs the SHACL shapes with RDFS entailment against isolated positive and intentionally invalid example ABoxes, and executes each competency-question query against the ontology plus positive examples. Complete expected answer rows make changed bindings, multiplicities, values, and unexamined empty results a validation failure even when the row count is unchanged. The same check runs in CI as the pull-request gate, and a tag-triggered workflow attaches the distributions to a GitHub release. This makes the resource's quality claims reproducible by any third party.

\textbf{Availability (FAIR).} HCMO targets the Resources-Track availability criteria \cite{fair}:

\begin{itemize}
\item \emph{Findable / Interoperable identification.} The ontology and its terms are
identified by a persistent \nolinkurl{w3id.org} IRI (\nolinkurl{https://w3id.org/hcmo/ontology/hcm#}, which dereferences via HTTP 303 content negotiation to the documentation); the resource is archived on Zenodo with a DOI.

\item \emph{Accessible.} Sources and distributions are openly available from the public
Git repository and the Zenodo deposit (download), with HTML documentation served via GitHub Pages. A public \textbf{SPARQL endpoint} will be provided for live querying. \emph{[pending: host the endpoint --- T6b.]}

\item \emph{Interoperable.} The current model reuses BFO \cite{bfo}, IAO, and SOSA/SSN
\cite{sosa} anchors, with Schema.org; example data reuse OWL-Time \cite{owltime}. SOSA is pinned to the 2017 Recommendation and uses \nolinkurl{sosa:ObservableProperty}. Canonical SemTS 1.2.0 segment and dimension terms are selectively reused for location result tables. A pinned example-level evidence slice uses PROV-O \cite{provo}, OBI, STATO, and ISA/Bioschemas without asserting HCMO class mappings or formal ISA RO-Crate conformance. The release also uses a pinned QUDT 3.4.0 quantity-value pattern for dimensions, sampling rates, specifications, and measured results. Broader vocabulary mappings remain future alignment work. A JSON-LD context is shipped for application developers exchanging data.

\item \emph{Reusable.} The resource is licensed CC BY 4.0, carries provenance on the
ontology header (creators with ORCIDs, version IRI), and provides a canonical citation (\nolinkurl{CITATION.cff}) plus a DOI.

\end{itemize}

Human- and machine-readable documentation is generated from the merged graph with WIDOCO \cite{widoco2017} (overview, term cross-reference, namespace declarations, a WebVOWL diagram, and provenance), published at the project's documentation site.

\textbf{Sustainability and governance.} HCMO is \textbf{lab-maintained} by the Huzard group: development is on GitHub with public issues and pull requests, versioning follows SemVer mirrored by \nolinkurl{owl:versionIRI}, and terms are deprecated rather than deleted or re-minted to protect external references. The COST \textbf{TEATIME} network serves as the community feedback channel grounding the model in real domain needs. A FAIR self-assessment (F/A/I/R) is given in the appendix (see \nolinkurl{metadata/resource-metadata.md}).

\begin{figure}[t]
\centering
\resizebox{\linewidth}{!}{%
\begin{tikzpicture}[
  node distance=4mm and 6mm,
  box/.style={draw, rounded corners=1.5pt, align=center, minimum height=8mm, font=\scriptsize, inner xsep=5pt},
  flow/.style={-{Latex[length=2mm]}, thick}
]
\node[box, fill=blue!8] (src) {Five source modules\\+ compatibility};
\node[box, fill=yellow!12, right=of src] (manifest) {Stable manifest\\\texttt{hcmo.yaml}};
\node[box, fill=green!10, right=of manifest] (build) {Deterministic build\\\texttt{build.py}};
\node[box, fill=green!10, right=of build] (dist) {TTL / OWL / JSON-LD\\profile inventory};
\node[box, fill=red!7, below=of build] (validate) {Parse + SHACL + 11 CQs\\audits + HermiT + ISA tests};
\node[box, fill=violet!8, right=of validate] (release) {CI gate\\tag + Zenodo + docs};
\draw[flow] (src) -- (manifest);
\draw[flow] (manifest) -- (build);
\draw[flow] (build) -- (dist);
\draw[flow] (dist) |- (validate);
\draw[flow] (validate) -- (release);
\draw[flow] (manifest) |- (validate);
\end{tikzpicture}
}
\caption{HCMO's manifest-driven, reproducible build and validation pipeline.}
\label{fig:pipeline}
\end{figure}

